% \iffalse meta-comment
%
% hit-fonts.dtx 是各个类共用的字体层
%
% \fi
%
% \section{共享字体层}
%
% \changes{v3.2a}{0000/00/00}{字体设置合成一份共享层，各个类共用}
% \changes{v3.2a}{0000/00/00}{新增 \option{font-en} 与 \option{font-zh}，西文与中文
%   字体可以分别指定，不再只能整体转给 \pkg{ctex}}
% \changes{v3.2a}{0000/00/00}{废弃模板自带的 \texttt{siyuan} fontset，改用
%   \option{font-zh} 取 \texttt{ubuntu}}
% 各个类共用的字体设置。西文与中文各由 \option{font-en} 与 \option{font-zh} 决定，
% 并提供了一套中文字号命令（字号没有统一规定，\hithesis{} 与 \CTeX、MS~Word、
% WPS、Adobe 的选择一致）。
%
% 本文件排在每个 \cs{file} 列表的第二位，此时 \pkg{fontspec} 与 \pkg{xeCJK} 尚未
% 加载，所以只做声明，要执行的放进 \cs{hit@setup@fonts}，由各个类调用。
%
%    \begin{macrocode}
%<*shared-fonts>
\bool_new:N \g__hit_platform_mac_bool
\str_new:N \l__hit_font_en_str
\str_new:N \l__hit_font_zh_str
\cs_new_protected:Npn \hit@setup@fonts
  {
%    \end{macrocode}
%
% 平台探测。\pkg{expl3} 只能判断是不是 Windows，macOS 看系统字体文件在不在。
%    \begin{macrocode}
\sys_if_platform_unix:T
  {
    \file_if_exist:nT { /System/Library/Fonts/Menlo.ttc }
      { \bool_gset_true:N \g__hit_platform_mac_bool }
  }
%    \end{macrocode}
%
% \changes{v3.2a}{0000/00/00}{\option{font-en} 取 \texttt{auto} 时改为除开源字体集外
%   一律优先 Times New Roman，原先 \texttt{adobe} 与 \texttt{founder} 拿不到}
% \option{font-en} 取 \texttt{auto} 时，\option{fontset} 是开源字体集
% （\texttt{fandol}、\texttt{ubuntu}）就用 TeX Gyre Termes，其余优先
% Times New Roman，找不到再退回。
%    \begin{macrocode}
\str_set:Nx \l__hit_font_en_str { \hit@font@en }
\str_if_eq:VnT \l__hit_font_en_str { } { \str_set:Nn \l__hit_font_en_str { auto } }
\str_if_eq:VnT \l__hit_font_en_str { auto }
  {
    \str_case:enF { \hit@fontset }
      {
        { fandol } { \str_set:Nn \l__hit_font_en_str { termes } }
        { ubuntu } { \str_set:Nn \l__hit_font_en_str { termes } }
      }
      { \str_set:Nn \l__hit_font_en_str { times } }
  }
\str_if_eq:VnT \l__hit_font_en_str { times }
  {
    \fontspec_font_if_exist:nF { Times~New~Roman }
      { \str_set:Nn \l__hit_font_en_str { termes } }
  }
\cs_new_protected:Npn \__hit_set_font_en_times:
  {
    \setmainfont { Times~New~Roman }
    \setsansfont { Arial }
    \setmonofont { Courier~New } [ Scale = MatchLowercase ]
  }
\cs_new_protected:Npn \__hit_set_font_en_termes:
  {
    \setmainfont { TeXGyreTermesX }
      [
        UprightFont    = *-Regular ,
        BoldFont       = *-Bold ,
        ItalicFont     = *-Italic ,
        BoldItalicFont = *-BoldItalic ,
        Extension      = .otf ,
      ]
    \setsansfont { texgyreheros }
      [
        UprightFont    = *-regular ,
        BoldFont       = *-bold ,
        ItalicFont     = *-italic ,
        BoldItalicFont = *-bolditalic ,
        Extension      = .otf ,
      ]
    \setmonofont { texgyrecursor }
      [
        UprightFont    = *-regular ,
        BoldFont       = *-bold ,
        ItalicFont     = *-italic ,
        BoldItalicFont = *-bolditalic ,
        Extension      = .otf ,
      ]
  }
\cs_new_protected:Npn \__hit_set_font_en_none: { }
%    \end{macrocode}
%
% \changes{v3.2a}{0000/00/00}{\option{font-en} 取值无效时报错，原先被 \cs{use:c}
%   静默吞掉，拼错的人拿到一份没设西文字体的文档}
% 取值无效时报错。
%    \begin{macrocode}
\cs_if_exist:cTF { __hit_set_font_en_ \l__hit_font_en_str : }
  { \use:c { __hit_set_font_en_ \l__hit_font_en_str : } }
  { \msg_error:nnx { hithesis } { font-en-unknown } { \l__hit_font_en_str } }
\cs_new_protected:Npn \__hit_set_font_zh_windows:
  {
    \setCJKmainfont { SimSun } [ AutoFakeBold = true , AutoFakeSlant = true ]
    \setCJKsansfont { SimHei } [ AutoFakeBold = true ]
    \setCJKmonofont { FangSong }
    \setCJKfamilyfont { zhsong } { SimSun } [ AutoFakeBold = true ]
    \setCJKfamilyfont { zhhei }  { SimHei } [ AutoFakeBold = true ]
    \setCJKfamilyfont { zhkai }  { KaiTi }
    \setCJKfamilyfont { zhfs }   { FangSong }
  }
\cs_new_protected:Npn \__hit_set_font_zh_mac:
  {
    \defaultCJKfontfeatures { }
    \setCJKmainfont { Songti~SC }
      [
        UprightFont    = *~Light ,
        BoldFont       = *~Bold ,
        AutoFakeSlant  = true ,
      ]
    \setCJKsansfont { Heiti~SC } [ UprightFont = *~Light , BoldFont = *~Medium ]
    \setCJKmonofont { STFangsong }
    \setCJKfamilyfont { zhsong } { Songti~SC }
      [ UprightFont = *~Light , BoldFont = *~Bold ]
    \setCJKfamilyfont { zhhei } { Heiti~SC }
      [ UprightFont = *~Light , BoldFont = *~Medium ]
    \setCJKfamilyfont { zhfs }  { STFangsong }
    \setCJKfamilyfont { zhkai } { Kaiti~SC }
      [ UprightFont = *~Regular , BoldFont = *~Bold ]
  }
\cs_new_protected:Npn \__hit_set_font_zh_ubuntu:
  {
    \defaultCJKfontfeatures { }
    \setCJKmainfont { Noto~Serif~CJK~SC }
      [
        UprightFont    = *~Light ,
        BoldFont       = *~Bold ,
        AutoFakeSlant  = true ,
        Script         = CJK ,
      ]
    \setCJKsansfont { Noto~Sans~CJK~SC } [ BoldFont = *~Medium , Script = CJK ]
    \setCJKmonofont { Noto~Sans~Mono~CJK~SC } [ Script = CJK ]
    \setCJKfamilyfont { zhsong } { Noto~Serif~CJK~SC }
      [ UprightFont = *~Light , BoldFont = *~Bold , Script = CJK ]
    \setCJKfamilyfont { zhhei } { Noto~Sans~CJK~SC }
      [ BoldFont = *~Medium , Script = CJK ]
    \fontspec_font_if_exist:nTF { LXGWZhuqueFangsong-Regular.ttf }
      {
        \setCJKmonofont { LXGWZhuqueFangsong-Regular.ttf }
        \setCJKfamilyfont { zhfs } { LXGWZhuqueFangsong-Regular.ttf }
      }
      {
        \setCJKfamilyfont { zhfs } { FandolFang } [ Extension = .otf , UprightFont = *-Regular ]
      }
    \fontspec_font_if_exist:nTF { AR~PL~KaitiM~GB }
      { \setCJKfamilyfont { zhkai } { AR~PL~KaitiM~GB } }
      { \setCJKfamilyfont { zhkai } { FandolKai } [ Extension = .otf , UprightFont = *-Regular ] }
  }
%    \end{macrocode}
%
% \changes{v3.2a}{0000/00/00}{新增 \texttt{hanyi} 一档，跟 \pkg{ctex} 2.6.4 的
%   同名 fontset 一致}
% 汉仪。\pkg{ctex} 2.6.4（2026）新加的一档。
%    \begin{macrocode}
\cs_new_protected:Npn \__hit_set_font_zh_hanyi:
  {
    \defaultCJKfontfeatures { }
    \setCJKmainfont { HYShuSongEr~S }
      [ BoldFont = HYZhongSong~S , AutoFakeSlant = true ]
    \setCJKsansfont { HYZhongHei~S } [ BoldFont = HYDaHei~S ]
    \setCJKmonofont { HYFangSong~S }
    \setCJKfamilyfont { zhsong } { HYShuSongEr~S } [ BoldFont = HYZhongSong~S ]
    \setCJKfamilyfont { zhhei }  { HYZhongHei~S }  [ BoldFont = HYDaHei~S ]
    \setCJKfamilyfont { zhkai }  { HYKaiTi~S }
    \setCJKfamilyfont { zhfs }   { HYFangSong~S }
  }
\cs_new_protected:Npn \__hit_set_font_zh_fandol:
  {
    \defaultCJKfontfeatures { }
    \setCJKmainfont { FandolSong }
      [
        Extension      = .otf ,
        UprightFont    = *-Regular ,
        BoldFont       = *-Bold ,
        AutoFakeSlant  = true ,
      ]
    \setCJKsansfont { FandolHei } [ Extension = .otf , UprightFont = *-Regular , BoldFont = *-Bold ]
    \setCJKmonofont { FandolFang } [ Extension = .otf , UprightFont = *-Regular ]
    \setCJKfamilyfont { zhsong } { FandolSong }
      [ Extension = .otf , UprightFont = *-Regular , BoldFont = *-Bold ]
    \setCJKfamilyfont { zhhei } { FandolHei }
      [ Extension = .otf , UprightFont = *-Regular , BoldFont = *-Bold ]
    \setCJKfamilyfont { zhfs }  { FandolFang } [ Extension = .otf , UprightFont = *-Regular ]
    \setCJKfamilyfont { zhkai } { FandolKai }  [ Extension = .otf , UprightFont = *-Regular ]
  }
%    \end{macrocode}
%
% Adobe 与方正两档。
%    \begin{macrocode}
\cs_new_protected:Npn \__hit_set_font_zh_adobe:
  {
    \defaultCJKfontfeatures { }
    \setCJKmainfont { AdobeSongStd-Light }
      [ BoldFont = AdobeHeitiStd-Regular , AutoFakeSlant = true ]
    \setCJKsansfont { AdobeHeitiStd-Regular }
    \setCJKmonofont { AdobeFangsongStd-Regular }
    \setCJKfamilyfont { zhsong } { AdobeSongStd-Light }
    \setCJKfamilyfont { zhhei }  { AdobeHeitiStd-Regular }
    \setCJKfamilyfont { zhfs }   { AdobeFangsongStd-Regular }
    \setCJKfamilyfont { zhkai }  { AdobeKaitiStd-Regular }
  }
\cs_new_protected:Npn \__hit_set_font_zh_founder:
  {
    \defaultCJKfontfeatures { }
    \setCJKmainfont { FZShuSong-Z01 }
      [ BoldFont = FZXiaoBiaoSong-B05 , AutoFakeSlant = true ]
    \setCJKsansfont { FZXiHeiI-Z08 } [ BoldFont = FZHei-B01 ]
    \setCJKmonofont { FZFangSong-Z02 }
    \setCJKfamilyfont { zhsong } { FZShuSong-Z01 } [ BoldFont = FZXiaoBiaoSong-B05 ]
    \setCJKfamilyfont { zhhei }  { FZHei-B01 }
    \setCJKfamilyfont { zhfs }   { FZFangSong-Z02 }
    \setCJKfamilyfont { zhkai }  { FZKai-Z03 }
  }
\cs_new_protected:Npn \__hit_set_font_zh_none: { }
%    \end{macrocode}
%
% \changes{v3.2a}{0000/00/00}{中文的斜体改用 \texttt{AutoFakeSlant}，不再拿楷体顶。
%   楷体仍旧由 \cs{kaishu} 与 \texttt{zhkai} 字族提供}
% \changes{v3.2a}{0000/00/00}{\texttt{ubuntu} 一档的仿宋与楷体跟 \pkg{ctex} 2.6.4
%   对齐：仿宋用霞鹜朱雀、楷体用 \texttt{AR~PL~KaitiM~GB}，都是装了才用，没装退回
%   Fandol。\pkg{ctex} 的同名 fontset 是硬指定，缺字体直接报错，这里不这么做}
% 中文字体。不写 \option{font-zh} 则整块跳过，字体由 \pkg{ctex} 按
% \option{fontset} 设置。取值一律沿用 \pkg{ctex} 的 fontset 名。
%
% 斜体走 \texttt{AutoFakeSlant}，倾斜因子取 \pkg{xeCJK} 的默认 0.167。
% \pkg{ctex} 与多数中文模板是把楷体指定成斜体，这里不这么做：楷体是独立的一族，
% 由 \cs{kaishu} 显式选用，不该被 \cs{textit} 顺带占掉。
%    \begin{macrocode}
\str_set:Nx \l__hit_font_zh_str { \hit@font@zh }
\str_if_eq:VnF \l__hit_font_zh_str { }
  {
    \str_if_eq:VnT \l__hit_font_zh_str { auto }
      {
        \str_case:enF { \hit@fontset }
          {
            { windows } { \str_set:Nn \l__hit_font_zh_str { windows } }
            { mac }     { \str_set:Nn \l__hit_font_zh_str { mac } }
            { ubuntu }  { \str_set:Nn \l__hit_font_zh_str { ubuntu } }
            { fandol }  { \str_set:Nn \l__hit_font_zh_str { fandol } }
            { adobe }   { \str_set:Nn \l__hit_font_zh_str { adobe } }
            { founder } { \str_set:Nn \l__hit_font_zh_str { founder } }
            { hanyi }   { \str_set:Nn \l__hit_font_zh_str { hanyi } }
          }
          {
            \sys_if_platform_windows:TF
              { \str_set:Nn \l__hit_font_zh_str { windows } }
              {
                \bool_if:NTF \g__hit_platform_mac_bool
                  { \str_set:Nn \l__hit_font_zh_str { mac } }
                  { \str_set:Nn \l__hit_font_zh_str { fandol } }
              }
          }
      }
    \cs_if_exist:cTF { __hit_set_font_zh_ \l__hit_font_zh_str : }
      { \use:c { __hit_set_font_zh_ \l__hit_font_zh_str : } }
      { \msg_error:nnx { hithesis } { font-zh-unknown } { \l__hit_font_zh_str } }
%    \end{macrocode}
%
% 字体由模板设置时传给 \pkg{ctex} 的是 \texttt{fontset=none}，这四个命令
% \pkg{ctex} 不会定义，要自己补。
%    \begin{macrocode}
    \cs_set_protected:Npn \songti { \CJKfamily { zhsong } }
    \cs_set_protected:Npn \heiti { \CJKfamily { zhhei } }
    \cs_set_protected:Npn \fangsong { \CJKfamily { zhfs } }
    \cs_set_protected:Npn \kaishu { \CJKfamily { zhkai } }
  }% \str_if_eq:VnF \l__hit_font_zh_str { } 到此为止
%    \end{macrocode}
%
% \changes{v3.2a}{0000/00/00}{把中文里常用的几个几何符号归成汉字，原先它们去找西文
%   字体，字体里没有就整个排丢了}
% \pkg{xeCJK} 按 Unicode 区块分中西文，「几何图形」这一块归了西文，于是中文正文里
% 写 \texttt{□}（规范原文用它占位，如「攻读□士学位」）会去找西文字体，
% \pkg{TeX~Gyre~Termes} 里没有这个字，整个字就排丢了——不是排成豆腐块，是什么都
% 没有，只在日志里留一句 Missing character。归成汉字之后从中文字体取，正常排出来。
%    \begin{macrocode}
    \clist_map_inline:nn { "25A1 , "25A0 , "25CB , "25CF , "25B3 , "25B2 }
      { \xeCJKDeclareCharClass { CJK } { ##1 } }
  }% \hit@setup@fonts 到此为止
%    \end{macrocode}
%
% MS~Word 里的 1\,pt 就是 \LaTeX{} 里的 1\,bp，而 1\,bp 等于 \LaTeX{} 里的
% 72.27/72\,pt。所以 Word 上看到的磅数，在这里直接写成 bp 即可。对照如下：
% \begin{center}
% \begin{tabular}{llll}
% \toprule
% 初号 & 42bp & 14.82mm & 42.1575pt \\
% 小初 & 36bp & 12.70mm & 36.135 pt \\
% 一号 & 26bp & 9.17mm & 26.0975pt \\
% 小一 & 24bp & 8.47mm & 24.09pt \\
% 二号 & 22bp & 7.76mm & 22.0825pt \\
% 小二 & 18bp & 6.35mm & 18.0675pt \\
% 三号 & 16bp & 5.64mm & 16.06pt \\
% 小三 & 15bp & 5.29mm & 15.05625pt \\
% 四号 & 14bp & 4.94mm & 14.0525pt \\
% 小四 & 12bp & 4.23mm & 12.045pt \\
% 五号 & 10.5bp & 3.70mm & 10.59375pt \\
% 小五 & 9bp & 3.18mm & 9.03375pt \\
% 六号 & 7.5bp & 2.56mm & \\
% 小六 & 6.5bp & 2.29mm & \\
% 七号 & 5.5bp & 1.94mm & \\
% 八号 & 5bp & 1.76mm & \\\bottomrule
% \end{tabular}
% \end{center}
%
% \changes{v3.2a}{0000/00/00}{\cs{hit@def@fontsize} 改名
%   \cs{\_\_hit\_define\_font\_size:nn}，它只在本文件里用}
% 定义一个字号命令，参数为名称与磅数。字号与行距解耦，行距由可选参数给出，
% 默认 1.3 倍。下面十七个字号的取值，\hithesis{} 与 \CTeX、MS~Word、WPS、Adobe
% 的选择一致，模板自己只用到其中十一个，其余留给用户。
%    \begin{macrocode}
\cs_new_protected:Npn \__hit_define_font_size:nn #1#2
  {
    \exp_args:Nc \NewDocumentCommand {#1} { O{1.3} }
      { \fontsize {#2} { ##1 \dimexpr #2 } \selectfont }
  }
%    \end{macrocode}
%
%
% \changes{v3.2a}{0000/00/00}{删掉 \cs{dachu}、\cs{banxiaosi}、\cs{dawu} 三个字号。
%   这三个不在通行的中文字号里，手册的字号对照表也没列，模板自己一处没用过}
% \begin{macro}{\chuhao}
% \begin{macro}{\xiaochu}
% \begin{macro}{\yihao}
% \begin{macro}{\xiaoyi}
% \begin{macro}{\erhao}
% \begin{macro}{\xiaoer}
% \begin{macro}{\sanhao}
% \begin{macro}{\xsanhao}
% \changes{v3.1b}{2022/06/06}{增加 \cs{xsanhao} 为了深圳校区的原创性声明标题}
% \begin{macro}{\xiaosan}
% \begin{macro}{\sihao}
% \begin{macro}{\xiaosi}
% \begin{macro}{\wuhao}
% \begin{macro}{\xiaowu}
% \begin{macro}{\liuhao}
% \begin{macro}{\xiaoliu}
% \begin{macro}{\qihao}
% \begin{macro}{\bahao}
% 一组字号定义。
%    \begin{macrocode}
\__hit_define_font_size:nn{chuhao}{42bp}
\__hit_define_font_size:nn{xiaochu}{36bp}
\__hit_define_font_size:nn{yihao}{26bp}
\__hit_define_font_size:nn{xiaoyi}{24bp}
\__hit_define_font_size:nn{erhao}{22bp}
\__hit_define_font_size:nn{xiaoer}{18bp}
\__hit_define_font_size:nn{sanhao}{16bp}
\__hit_define_font_size:nn{xsanhao}{15.9bp}
\__hit_define_font_size:nn{xiaosan}{15bp}
\__hit_define_font_size:nn{sihao}{14bp}
\__hit_define_font_size:nn{xiaosi}{12bp}
\__hit_define_font_size:nn{wuhao}{10.5bp}
\__hit_define_font_size:nn{xiaowu}{9bp}
\__hit_define_font_size:nn{liuhao}{7.5bp}
\__hit_define_font_size:nn{xiaoliu}{6.5bp}
\__hit_define_font_size:nn{qihao}{5.5bp}
\__hit_define_font_size:nn{bahao}{5bp}
%    \end{macrocode}
%
% \cs{xsanhao} 比三号略小，给深圳校区本科生的原创性声明标题用，取自
% \hitszthesis{}。规范里那个标题是三号，但 \cs{sanhao} 一行排不下。
%    \begin{macrocode}
%</shared-fonts>
%    \end{macrocode}
%
% \end{macro}
% \end{macro}
% \end{macro}
% \end{macro}
% \end{macro}
% \end{macro}
% \end{macro}
% \end{macro}
% \end{macro}
% \end{macro}
% \end{macro}
% \end{macro}
% \end{macro}
% \end{macro}
% \end{macro}
% \end{macro}
% \end{macro}
%
%
% \Finale
